\section{Short-shift radial energy and spectral laws}
\label{sec:short-shift-spectrum}

The benefit of the added coordinate appears before any estimate is used.

\begin{proposition}[Exact product-center identity]
\label{prop:exact-product-center-identity}
For every $h,D,X$,
\begin{equation}
 \cF_{h,D}(X)
 =\sum_rA_D(r)(\Lambda(r+h)-1)^2.
 \label{eq:radial-energy-product-center}
\end{equation}
If
\begin{equation}
 \cF^{\raw}_{h,D}(X):=
 \sum_{\substack{n>D\\m\geq1}}
 u(n)^2\Lambda(mn+h)^2|W(mn/X)|^2,
 \label{eq:raw-radial-energy}
\end{equation}
then
\begin{equation}
 \cF^{\raw}_{h,D}(X)=\sum_rA_D(r)\Lambda(r+h)^2.
 \label{eq:raw-radial-product-center}
\end{equation}
\end{proposition}

\begin{proof}
By \cref{prop:coordinate-bijection}, summing over primitive triples is
the same as summing once over every factor pair $(m,n)$.  Group those
pairs by $r=mn$.  The source squares on one product center sum to
$A_D(r)$, proving both identities.
\end{proof}

No target with a different radial index occurs in
\eqref{eq:radial-energy-product-center}.  In particular, no two-target
upper sieve is needed.

\begin{theorem}[Short-shift fully radial energy]
\label{thm:short-shift-radial-energy}
Let $G\in C_c^1((0,\infty))$ be nonnegative and nonzero.  Fix
$0<\eps<13/15$ and $0<\eta<1$.  Uniformly for
\eqref{eq:intro-parameter-range},
\begin{equation}
 \begin{aligned}
 \sum_hG(h/H)\cF_{h,D}(X)
 ={}&\frac12HX I_0(G)I_2(W)\log X\\
 &\times\bigl((\log X)^2-(\log D)^2\bigr)
 +O_{\eps,\eta,W,G}\bigl(HX(\log X)^2\bigr).
 \end{aligned}
 \label{eq:short-shift-radial-energy}
\end{equation}
The same asymptotic holds for $\cF^{\raw}_{h,D}(X)$.
\end{theorem}

\begin{proof}
Insert \eqref{eq:radial-energy-product-center} and apply
\cref{prop:weighted-target-square}.  This gives
\[
 \sum_hG(h/H)\cF_{h,D}(X)
 =HI_0(G)\log X\,\cS_D(X)
 +O(HX(\log X)^2).
\]
Now substitute \eqref{eq:source-diagonal}.  Its remainder, after
multiplication by $H\log X$, is also $O(HX(\log X)^2)$.  The raw result
follows from \eqref{eq:raw-radial-product-center} and the raw part of
\cref{prop:weighted-target-square}.
\end{proof}

If $W\not\equiv0$, then $I_2(W)>0$ and
\[
 (\log X)^2-(\log D)^2
 \geq(2\eta-\eta^2)(\log X)^2.
\]
Thus the relative error in \eqref{eq:short-shift-radial-energy} is
$O_{\eta}(1/\log X)$.

\begin{remark}[A changed observable, not a stronger estimate of the old one]
The ratio-only ray energy of TPC-12 is
\begin{equation}
 \cE_{h,D}(X)
 =\sum_{(a,b)=1}\left|\sum_kx_{a,b,k}\right|^2
 =\cF_{h,D}(X)
 +\sum_{(a,b)=1}\sum_{k\ne\ell}
 x_{a,b,k}\overline{x_{a,b,\ell}}.
 \label{eq:old-versus-new-energy}
\end{equation}
The $O(HX(\log X)^2\log\log X)$ error in TPC-12 bounds the last sum
after shift averaging by a two-linear-forms Selberg upper sieve.  The
present $O(HX(\log X)^2)$ error is smaller because
\eqref{eq:intro-radial-energy} defines a different, fully separated
observable in which the off-diagonal is absent.  We do not claim an
improved error for $\cE_{h,D}$ itself.
\end{remark}

\subsection{Global two-dimensional Fej\'er band}

Combining the arithmetic theorem with
\cref{cor:global-radial-Parseval} gives a single-band statement.

\begin{corollary}[Short-shift global joint band]
\label{cor:short-shift-global-band}
If $T\delta_X\geq1$ and $S\Delta_X\geq1$, then under the hypotheses of
\cref{thm:short-shift-radial-energy},
\begin{equation}
 \begin{aligned}
 &\sum_hG(h/H)\iint_{\R^2}
 \kappa(t/T)\kappa(s/S)
 |\cR^{\rad}_{h,D}(X;t,s)|^2\,dt\,ds\\
 &\quad={}
 \frac12TSHX I_0(G)I_2(W)\log X
 \bigl((\log X)^2-(\log D)^2\bigr)\\
 &\qquad\quad+O_{\eps,\eta,W,G}
 \bigl(TSHX(\log X)^2\bigr).
 \end{aligned}
 \label{eq:short-shift-global-band}
\end{equation}
The same statement holds for the raw transform.
\end{corollary}

\begin{proof}
For every $h$, \eqref{eq:global-radial-Parseval} turns the integral into
$TS\cF_{h,D}(X)$.  Sum \cref{thm:short-shift-radial-energy}.
\end{proof}

\subsection{Sector-adapted multiresolution band}

For each nonempty $R$-sector, choose $T_R,S_R$ satisfying
\eqref{eq:two-thresholds}.  Define
\begin{equation}
 \begin{aligned}
 \mathfrak P_{h,D}(X)
 :={}&\sum_{\substack{R\in\cD_X\\\mathscr I_R(X)\ne\varnothing}}
 \frac1{T_RS_R}
 \iint_{\R^2}\kappa(t/T_R)\kappa(s/S_R)\\
 &\hspace{6em}\times
 |\cR^{\rad}_{h,D;R}(X;t,s)|^2\,dt\,ds.
 \end{aligned}
 \label{eq:multiresolution-functional}
\end{equation}
Equation \eqref{eq:multiresolution-Parseval} says exactly that
\begin{equation}
 \mathfrak P_{h,D}(X)=\cF_{h,D}(X).
 \label{eq:multiresolution-equals-energy}
\end{equation}
We therefore obtain the sector-optimal version without needing an
individual arithmetic asymptotic for each $R$.

\begin{corollary}[Short-shift multiresolution spectrum]
\label{cor:short-shift-multiresolution}
Under \eqref{eq:intro-parameter-range},
\begin{equation}
 \begin{aligned}
 \sum_hG(h/H)\mathfrak P_{h,D}(X)
 ={}&\frac12HX I_0(G)I_2(W)\log X\\
 &\times\bigl((\log X)^2-(\log D)^2\bigr)
 +O\bigl(HX(\log X)^2\bigr).
 \end{aligned}
 \label{eq:short-shift-multiresolution}
\end{equation}
\end{corollary}

This corollary combines sectorwise spectral resolutions only after each
sector has been isolated and normalized.  It neither asserts a common
small global band nor evaluates a specified fixed shift.
